% ============================ 5. GAPS ==================================
\newpage
\phase{Section five}{What India is missing}{Specific needs in the Indian market that nobody serves well yet, each with the evidence that it is real.}

{\tablefont
\begin{xltabular}{\textwidth}{@{}>{\RaggedRight\arraybackslash}p{24mm} L L c c@{}}
\caption{Unmet needs in the Indian green hydrogen market. ``Size'' is how big the need is; ``Fit'' is how well a capital-light start-up could serve it (both 1–5).}\\
\toprule
\thead{Gap} & \thead{Evidence it is real} & \thead{Who would pay} & \thead{Size} & \thead{Fit}\\ \midrule
\endfirsthead
\toprule
\thead{Gap} & \thead{Evidence it is real} & \thead{Who would pay} & \thead{Size} & \thead{Fit}\\ \midrule
\endhead
\textbf{1. Building what has been allocated} & 862,000 t/yr allocated but about 8,000 t/yr commissioned by Feb 2026 \cite{sqnghm2025,chemindigest2026} \rep; AM Green's first exports moved to 2028 \cite{amgreenuniper2026} & Developers, lenders & 5 & 2\\
\textbf{2. Stack servicing and independent testing} & Eight Indian stack makers won incentives, some at only 100 MW a year \cite{secipli2024}; global OEMs are under strain (Cummins exit \cite{argus2026}, Enapter returns \cite{enapter2026}); only 5 testing facilities approved nationally \cite{sqnghm2025} \rep. A plant whose OEM fails has nobody to service it & Plant owners, lenders, insurers & 4 & 4\\
\textbf{3. Certification and data} & India's certification scheme launched in 2025, run by BEE through accredited verifiers \cite{ghci2025}; EU buyers such as Uniper want RFNBO-certified ammonia \cite{amgreenuniper2026}; the EU border carbon charge covers hydrogen and ammonia \cite{coolset}. Exporters need to satisfy two rulebooks with hourly data & Producers, exporters & 3 & 5\\
\textbf{4. Mid-sized industrial users} & Glass, electronics, chemicals and chlor-alkali users are many and small \cite{igef2021}; Jindal Stainless already buys green hydrogen from an on-site plant built, owned and operated by Hygenco \cite{hygenco2024} \rep & Steel, glass, food, pharma, electronics plants & 3 & 4\\
\textbf{5. Electrolyser-ready water} & 20–30 L of water per kg once purification is counted \cite{rmi2023}; port hubs at Kandla, Paradip and Tuticorin are coastal \cite{rw2025} & Developers, EPCs & 3 & 3\\
\textbf{6. Demand certainty} & Refinery tenders were slow to arrive in 2025 \cite{bs2025}; purchase mandates are discussed but not notified \cite{xynteo2026} & The state & 5 & 1\\
\textbf{7. Small-ticket finance} & Big projects draw on REC (₹4,000 crore to ACME for Oman) \cite{acmeodisha2023} and the World Bank \cite{wb2023}; nothing comparable is visible for ₹10–100 crore on-site plants \estd & Banks, NBFCs, insurers & 3 & 2\\
\bottomrule
\end{xltabular}}

\begin{finding}[Key finding \textbullet\ the biggest gaps are not the ones a start-up should chase]
Gaps 1 and 6, building the allocated plants and creating demand, are the largest, but they need state power or billions of rupees. Gaps 2, 3 and 4 are smaller but match a capital-light, analytics-heavy team. They also reinforce each other. A company that services stacks collects the operating data; that data feeds certification; and certified, well-serviced plants are what a hydrogen-as-a-service contract sells. \confM
\end{finding}

\begin{plainwords}[An honest limit of this section]
The evidence for gaps 4 and 7 is thinner than for the others. I found one strong example of on-site green hydrogen for a mid-sized user (Jindal and Hygenco) and no good public data on what small Indian users pay for cylinder hydrogen today. Before building a business on gap 4, we would need to interview a dozen glass, food and electronics plants about their volumes and prices. That is a perfect field exercise for the research paper.
\end{plainwords}
